% =====================================================================
%  Conferenze — Evento 1: GAMM di Roma, Checkpoint Festival (3a ed.)
%  Introduzione all'evento e mappa dei capitoli C1.x.
%  Capitolo non numerato. Nessun pacchetto aggiuntivo.
% =====================================================================

\chapter*{GAMM di Roma: Checkpoint Festival, terza edizione}
\phantomsection
\addcontentsline{toc}{chapter}{GAMM di Roma: Checkpoint Festival, terza edizione}

Il \textbf{GAMM} (Game Museum) di Roma è un museo dedicato alla storia e alla
cultura del videogioco. Durante la \textbf{terza edizione del Checkpoint Festival}
ha ospitato una giornata di incontri che, a differenza della conferenza di Zagarolo
(capitoli C2.x e C3.x), guarda al videogioco soprattutto come \textbf{fenomeno
culturale e sociale}: tutela dei minori, storia, preservazione, musica, ma anche
l'esperienza concreta di chi sviluppa giochi in Italia.

Gli incontri registrati sono cinque:
\begin{enumerate}
  \item un panel sul rapporto tra videogiochi, società e \textbf{tutela dei minori},
        con istituzioni e industria (capitolo~\ref{cap:c1-1-tutela-minori});
  \item il percorso \textbf{dalla tesi di laurea al videogioco storico}
        (capitolo~\ref{cap:c1-2-videogioco-storico});
  \item dieci anni di \textbf{horror indipendente} raccontati dai fondatori di uno
        studio italiano (capitolo~\ref{cap:c1-3-horror-indipendente});
  \item un panel, in gran parte in inglese, sulla \textbf{preservazione del
        videogioco}: archivi, musei e distribuzione digitale
        (capitolo~\ref{cap:c1-4-preservazione}). La registrazione è divisa in due
        file, ma si tratta di un unico incontro;
  \item \textbf{musica e videogiochi}, dalla formazione al lavoro, che chiude il
        festival (capitolo~\ref{cap:c1-5-musica}).
\end{enumerate}
